\section*{الفصل \ref{ch:b2:wilson-cycle} --- الكتل القديمة ودورة ويلسون}

\begin{solution}{exo:b2:wilson-cycle:1}
(أ) $600 - 400 = \qty{200}{Myr}$، من انفصال الهوامش إلى آخر تصادم.
(ب) $480 - 400 = \qty{80}{Myr}$. (ج) التكوّن الجبلي الكاليدوني مجموعة من المراحل الموزّعة على
نحو 200 مليون سنة، وكل مرحلة بدأت في أوقات مختلفة في أماكن مختلفة: فالكلمة
تحدّد موقع محيط وحزامه، لا لحظة. والعمر يُعطى بالسلّم الطبقي أو
بملايين السنين.
\end{solution}

\begin{solution}{exo:b2:wilson-cycle:2}
(ج) شقّ شرق إفريقيا، التصدّع الخسفي؛ (أ) البحر الأحمر، محيط فتيّ؛ (د) الأطلسي، محيط
واسع بهوامش سلبية؛ (ه) الهامش الهادئي لأمريكا الجنوبية، الاندساس
وانغلاق محيط؛ (ب) الهيمالايا، التصادم؛ (و) بريتاني، حزام مبريّ، قد يبدأ بعده
التصدّع الخسفي من جديد.
\end{solution}

\begin{solution}{exo:b2:wilson-cycle:3}
(أ) بعد ترسّب طبقات أوائل السيلوري، أي بعد \qty{443.1}{Ma}، وقبل
ترسّب طبقات الديفوني المتأخر، أي قبل \qty{358.9}{Ma}، نهاية الديفوني.
(ب) أقل من $443.1 - 358.9 = \qty{84.2}{Myr}$؛ بل أقل في الواقع، لأن الطبقات السفلى أحدث
من قاعدة السيلوري والطبقات العليا أقدم من قمّته. (ج) ترسّب
التوربيديتات الرمادية في المحيط؛ فالطيّ والانفصام في التصادمات الكاليدونية؛ فالرفع
والحتّ حتى سطح أرضي غير مستوٍ؛ فترسّب البريشيا والكونغلوميرا الحمراء؛ فالإمالة
الخفيفة في التكوّن الجبلي الفاريسكي؛ فحتّ اليوم، الذي يقطع الجرف.
\end{solution}

\begin{solution}{exo:b2:wilson-cycle:4}
(أ) كلاهما: كتلة من القشرة القارية تفصلها عن القارات محيطات حين كانت
تعبر إيابيتوس، وهي بعد أن رست كتلة تحدّها صدوع ويختلف تاريخها عن تاريخ
جاراتها. (ب) لا هذا ولا ذاك: قارة، لا كتلة أُضيفت إلى قارة. (ج) \omterm{def:b2:wilson-cycle:terrane}{رقعة قشرية}، لا
\omterm{def:b2:wilson-cycle:terrane}{قارة صغرى}: قشرتها قشرة قوس، لا قطعة من قارة. (د) لا هذا ولا ذاك: إنه
الحدّ بين \omterm{def:b2:wilson-cycle:terrane}{الرقع القشرية}، أثر المحيط المنغلق.
\end{solution}

\begin{solution}{exo:b2:wilson-cycle:5}
(أ) $3000 - 2600 = 400$، $2600 - 1500 = 1100$، $1500 - 875 = 625$، $875 - 600 = 275$، $600 -
262.5 = 337.5$؛ والمتوسط $2737.5/5 = \qty{547.5}{Myr}$. (ب) $1100/275 = 4.0$. (ج) $(1500 -
262.5)/3 = \qty{412.5}{Myr}$. يتغيّر المتوسط بمقدار الربع حين تُستبعد النقاط المشكوك فيها،
وتتفاوت الفواصل بعامل أربعة: فهو تكرار تقريبي، لا دور.
\end{solution}

\begin{solution}{exo:b2:wilson-cycle:6}
(أ) $3000/(70 - 25) = \qty{66.7}{Myr}$. (ب) $3000/70 = \qty{42.9}{Myr}$. (ج) بعد أن تندسّ
الظهرة، يكون قاع المحيط الباقي جزءًا من الصفيحة البعيدة، التي تحمل القارة
البعيدة: فمعدّل الانغلاق هو سرعة تلك الصفيحة بالنسبة إلى الخندق.
\end{solution}

\begin{solution}{exo:b2:wilson-cycle:7}
(أ) $700/8 = \qty{87.5}{km}$؛ $700/30 = \qty{23.3}{km}$. (ب) الحزام الفاريسكي: فعند
\qty{30}{\celsius} لكل كيلومتر تنصهر القشرة الوسطى؛ وعند \qty{8}{\celsius} لكل كيلومتر
تقع حرارة الانصهار تحت قاعدة القشرة حتى لو كانت متثخّنة. (ج) يصف التدرّج المنخفض
صخورًا حُملت إلى الأسفل بسرعة في نطاق اندساس، قبل أن تتمكّن من التسخّن. أمّا الصخور التي
تبقى عميقة في قشرة متثخّنة فتدفأ نحو تدرّج حراري عادي على مدى ملايين السنين، ويمكنها
عندئذ أن تبلغ الانصهار على عمق أصغر.
\end{solution}

\begin{solution}{exo:b2:wilson-cycle:8}
(أ) $\tau\ln(3/0.3) = 26.3\times2.303 = \qty{60.6}{Myr}$. (ب) $290 - 60.6 = \qty{229}{Ma}$، في
الترياسي المتأخر. (ج) $\tau = (290 - 251.9)/\ln10 = \qty{16.5}{Myr}$. (د) الانهيار
التمدّدي للقشرة المتثخّنة، الذي يرقّقها ويخفض سطحها مباشرة؛ وكذلك
معدّل تعرية أعلى، أو مناخ أرطب، أو صخور ألين، أو ارتفاع ابتدائي أدنى.
\end{solution}

\begin{solution}{exo:b2:wilson-cycle:9}
(أ) الدرز، أثر إيابيتوس، يقع على اليابسة في بريطانيا وأيرلندا، شرق
الأطلسي؛ وقشرة لورنتيا، من الجهة الأخرى لإيابيتوس، تقع قرب روكول في
الجهة الأوروبية. فقد قطع الشقّ الجديد غرب الدرز القديم: انفتح قربه، لا على امتداده.
(ب) شمال الدرز، القاعدة والغطاء الرسوبي للهامش اللورنتي؛ وجنوبه، قاعدة
أفالونيا وغطاؤها، بتاريخين مختلفين، وبحيوانات مختلفة في أوائل الباليوزوي. (ج) جاء التصدّع الخسفي
بعد 80 مليون سنة، أي أكثر من 30 إلى 50 التي قد يصير بعدها النطاق البارد
قويًا؛ ومع ذلك تبع وسط الأطلسي جبال الأبلاش فعلًا. فحرارة الحزام تخبو؛ أمّا
صدوعه ونسيجه وتراكيبه فتبقى ويمكنها أن توجّه شقًّا.
\end{solution}

\begin{solution}{exo:b2:wilson-cycle:10}
(أ) $70 - 40/2 = \qty{50}{mm}$ في السنة. (ب) $\qty{1500}{km}/\qty{50}{km.Myr^{-1}} =
\qty{30}{Myr}$. (ج) يضيق المحيط بمعدّل $70 - 40 = \qty{30}{mm}$ في السنة: $3000 - 30\times30 =
\qty{2100}{km}$. (د) يصل عندئذ إلى الخندق غلاف صخري فتيّ حارّ، ولا تعود كتلة قديمة باردة
تفقد ماءها تحت القوس: فتتغيّر صهارية القوس أو تتوقّف.
\end{solution}

\begin{solution}{exo:b2:wilson-cycle:11}
(أ) $h = \Delta T(\rho_m - \rho_c)/\rho_m = 25\times600/3300 = \qty{4.5}{km}$. (ب) السمك الإضافي كله،
\qty{25}{km}، أي 5.5 أضعاف ارتفاع السلسلة. (ج) $2700\times9.8\times25\,000 =
\qty{6.6e8}{Pa}$، أي نحو \qty{0.66}{GPa}. (د) عند نحو \qty{0.7}{GPa} وحرارة
جذر قشرة متثخّنة، بضع مئات من الدرجات، سحنة الأمفيبوليت أو
الغرانوليت: نيسات الكتل القديمة وميغماتيتاتها.
\end{solution}

\begin{solution}{exo:b2:wilson-cycle:12}
(أ) $S = 30\times10 = \qty{300}{km}$؛ $W_0 = 300 + 300 = \qty{600}{km}$؛ $T = 35\times600/300 =
\qty{70}{km}$. (ب) $h = 35\times600/3300 = \qty{6.36}{km}$؛ $r = 4.5h = \qty{28.6}{km}$.
(ج) $\Delta F = 2700\times9.8\times6364\times(35\,000 + 3182 + 14\,318) =
\qty{8.8e12}{N.m^{-1}}$، أي 1.4 ضعف فائض التبت. (د) حزام كهذا أثقل من التبت، التي
تنبسط أصلًا: انهيار ثقالي، بتمدّد مناطقه العالية واندفاع عند حوافه.
\end{solution}

\begin{solution}{pb:b2:wilson-cycle:1}
\textbf{1.} $720/(540 - 500) = \qty{18}{km}$ في المليون سنة، أي \qty{18}{mm} في السنة.
\textbf{2.} \qty{36}{mm} في السنة: ظهرة بطيئة التوسّع (من 20 إلى \qty{55}{mm} في السنة).
\textbf{3.} $36\times(580 - 540) = \qty{1440}{km}$.
\textbf{4.} $36\times(580 - 500) = \qty{2880}{km}$.
\textbf{5.} توسّع متناظر بمعدّل ثابت من عرض صفري عند \qty{580}{Ma}، ولا
اندساس قبل \qty{500}{Ma}.
\textbf{6.} \qty{80}{Myr}، عند حافتي المحيط كلتيهما، ملاصقًا للهوامش القارية. وهو
أبرد غلاف صخري محيطي في المحيط وأسمكه وأكثفه، والأكثر استعدادًا للغوص.
\textbf{7.} $64 - 36 = \qty{28}{mm}$ في السنة.
\textbf{8.} $2880/28 = \qty{102.9}{Myr}$ بعد \qty{500}{Ma}: أي عند \qty{397.1}{Ma}.
\textbf{9.} الظهرة، على بعد \qty{1440}{km} من الخندق، تقترب منه بمعدّل $64 - 18 =
\qty{46}{mm}$ في السنة: بعد $1440/46 = \qty{31.3}{Myr}$، أي عند \qty{468.7}{Ma}.
\textbf{10.} $2880 - 28\times31.3 = \qty{2004}{km}$.
\textbf{11.} $2004/40 = \qty{50.1}{Myr}$ بعد ذلك: أي عند $468.7 - 50.1 = \qty{418.6}{Ma}$.
\textbf{12.} داخل المدى من 425 إلى \qty{400}{Ma} للتصادم الإسكاندي.
\textbf{13.} $40\times8 = \qty{320}{km}$.
\textbf{14.} $W_0 = 400 + 320 = \qty{720}{km}$؛ $T = 35\times720/400 = \qty{63.0}{km}$.
\textbf{15.} $h = 28\times600/3300 = \qty{5.09}{km}$؛ $r = 4.5\times5.09 = \qty{22.9}{km}$.
\textbf{16.} $\Delta F = 2700\times9.8\times5091\times(35\,000 + 2545 + 11\,455) =
\qty{6.60e12}{N.m^{-1}}$، أي أكثر من فائض التبت بنحو \qty{8}{\percent}.
\textbf{17.} أن يتوقّف عن الارتفاع وينبسط: تمدّد وصدوع عادية في داخله العالي،
وصدوع دفعية عند حوافه، كما تفعل التبت اليوم.
\textbf{18.} $f = 600/3300 = 0.182$؛ $k = 2.303/11 = 0.209$ لكل مليون سنة؛ $\tau = 1/(kf)
= \qty{26.3}{Myr}$.
\textbf{19.} $3.0\,e^{-30/26.3} = \qty{0.96}{km}$.
\textbf{20.} $26.3\ln(3.0/0.5) = \qty{47.1}{Myr}$ بعد \qty{410}{Ma}: أي عند \qty{362.9}{Ma}.
\textbf{21.} القشرة الإضافية هي $h/f$: إذ تهبط $3.0/0.182 = \qty{16.5}{km}$ إلى $0.5/0.182 =
\qty{2.75}{km}$، فيُزال \qty{13.75}{km} من الصخر.
\textbf{22.} $2700\times9.8\times13\,750 = \qty{3.6e8}{Pa}$، أي نحو \qty{0.36}{GPa}.
\textbf{23.} نعم: بحلول الديفوني المتأخر يكون متوسط ارتفاع الحزام النموذجي قد هبط إلى نحو \qty{1}{km}
ولا يزال يهبط، مع تضاريس باقية بين أوديته وحيوده؛ ويمكن للحصى الحمراء أن
تتجمّع في منخفضاته قبل \qty{358.9}{Ma} بزمن طويل، كما في سيكار بوينت، حيث ترتكز على سطح
ذي تضاريس.
\textbf{24.} \textbf{\qty[text-series-to-math=true]{2880}{km}}، عند \qty{500}{Ma}، حين
بدأ الاندساس.
\end{solution}
